\clearpage
\SourceLeaf{454}%
\pdfbookmark[0]{Gauss to Enke}{gauss-an-enke-1849}
\renewcommand{\BookRunningTitle}{Gauss to Enke}
\begingroup
\begin{center}{\large GAUSS TO ENKE.}\end{center}
\BookRule{.20\textwidth}
\begin{center}Highly esteemed friend!\end{center}

— — The remarks on the frequency of prime numbers that you kindly sent me interested me in more than one way. They reminded me of my own work on this subject, whose beginnings lie far in the past, in 1792 or 1793, when I acquired \textsc{Lambert}’s supplements to the logarithm tables. Even before I had begun the finer investigations of higher arithmetic, one of my first tasks was to turn my attention to the decreasing frequency of prime numbers. For this purpose I counted them in each group of a thousand and wrote the results on one of the attached blank leaves. I soon saw that, despite all its fluctuations, this frequency was on average approximately inversely proportional to the logarithm, so that the number of all primes below a given bound $n$ would be approximately expressed by the integral
\[
\int\frac{\mathrm{d}n}{\log n}
\]
with the natural logarithm understood. Later, when I learned of the list printed in \textsc{Vega}’s tables (of 1796), extending up to 400031, I extended my count, which confirmed that relationship. In 1811 the appearance of \textsc{Chernac}’s
\SourceLeaf{455}%
cribrum gave me great pleasure, and (since I had no patience for a sustained count in sequence) I very often spent an otherwise free quarter of an hour counting a group of a thousand, here or there; eventually, however, I gave it up altogether without quite finishing the million. Only later did I make use of \textsc{Goldschmidt}’s diligence, partly to fill the remaining gaps in the first million and partly to continue the count using \textsc{Burck\-hardt}’s tables Thus the first three millions have been counted (for many years now) and compared with the value of the integral. Here I give only a short extract:

\begin{center}
\begingroup\fontsize{10}{11.5}\selectfont\setlength{\tabcolsep}{3pt}
\begin{tabular}{r|r|rr|rr}
Below & \shortstack{there are\\prime numbers} & \shortstack{Integral\\$\displaystyle\int\frac{\mathrm{d}n}{\log n}$} & Deviation & \shortstack{Your\\formula} & Deviation \\
500000 & 41556 & $41606{,}4$ & $+50{,}4$ & $41596{,}9$ & $+40{,}9$ \\
1000000 & 78501 & $79627{,}5$ & $+126{,}5$ & $78672{,}7$ & $+171{,}7$ \\
1500000 & 114112 & $114263{,}1$ & $+151{,}1$ & $114374{,}0$ & $+264{,}0$ \\
2000000 & 148883 & $149054{,}8$ & $+171{,}8$ & $149233{,}0$ & $+350{,}0$ \\
2500000 & 183016 & $183245{,}0$ & $+229{,}0$ & $183495{,}1$ & $+479{,}1$ \\
3000000 & 216745 & $216970{,}6$ & $+225{,}6$ & $217308{,}5$ & $+563{,}5$
\end{tabular}\endgroup
\end{center}

\Needspace{9\baselineskip}
I did not know that \textsc{Legendre} had also concerned himself with this subject. Your letter prompted me to look in his \textit{Théorie des Nombres}, and in the second edition I found some pages on it that I must have overlooked earlier (or forgotten since). \textsc{Legendre} uses the formula
\[
\frac{n}{\log n-A}
\]
where $A$ is supposed to be a constant, for which he takes $1{,}08366$. A quick calculation then gives me the deviations in the above cases as
\[
\begin{array}{r}
-23{,}3\\
+42{,}2\\
+68{,}1\\
+92{,}8\\
+159{,}1\\
+167{,}6
\end{array}
\]

These differences are even smaller than those with the integral, but they seem to grow faster as $n$ increases, so it is quite possible
\SourceLeaf{456}%
that, if continued much further, they would exceed the latter. To bring the count and the formula into agreement, one would have to set, respectively, in place of \mbox{$A=1{,}08366$}
\[
\begin{array}{r}
1{,}09040\\
1{,}07682\\
1{,}07582\\
1{,}07529\\
1{,}07179\\
1{,}07297
\end{array}
\]

It seems that, as $n$ grows, the (average) value of $A$ decreases, but I venture no conjecture as to whether the limit as $n$ grows to infinity will be $1$ or a quantity different from $1$. I cannot say that there are grounds for expecting a very simple limiting value; on the other hand, the excess of $A$ over $1$ could quite well be a quantity of order $\frac{1}{\log n}$. I would be inclined to believe that the differential of the function in question must be simpler than the function itself. Whereas I have assumed $\frac{\mathrm{d}n}{\log n}$ for it, \textsc{Legendre}’s formula would assume a differential function that would be approximately $\frac{\mathrm{d}n}{\log n-(A-1)}$. Your formula, moreover, could for a very large $n$ be regarded as agreeing with
\[
\frac{n}{\log n-\frac{1}{2k}}
\]
where $k$ is the modulus of \textsc{Briggs}’s logarithms, and hence with \textsc{Legendre}’s formula if
\[
A=\tfrac{1}{2k}=1{,}1513\ \text{is set.}
\]

\Needspace{5\baselineskip}
Finally, I should mention that I have noticed a couple of differences between your counts and mine.
\begin{center}
\begin{tabular}{lrlrlrlr}
Between & 59000 & and & 60000 & you have & 95 & I & 94\\
 & 101000 & & 102000 & & 94 & & 93
\end{tabular}
\end{center}
The first difference may be due to the prime number 59023 appearing twice in \textsc{Lambert}’s supplements. The group of a thousand from 101000—102000 teems with errors in \textsc{Lambert}’s supplements: in my copy I have marked 7 numbers that are not prime and inserted 2 that were missing.
\SourceLeaf{457}%
Could you not persuade the young \textsc{Dase} to count the primes in the following millions from the tables at the Academy, which, I fear, the public is not meant to possess? In that case I should mention that, in the 2nd and 3rd millions, the count was made at my direction according to a special scheme that I had already used myself for part of the first million. The counts for each 100000 occupy one (small) octavo page in 10 columns, each referring to one group of ten thousand; there is also a column before them (on the left) and another after them on the right. As an example, here is one vertical column and the two additional columns from the interval $1000000\mathinner{.\,.\,.}1100000$ \mbox{— — —}

For an explanation, take the 1st vertical column, for example. The group of ten thousand from 1000000 to 1010000 contains 100 groups of a hundred; among these there is 1 containing only one prime; none at all with 2 or 3; 2 with 4 primes each; 11 with 5 each, and so on. Together they give $752=1\pdot1+4\pdot2+5\pdot11+6\pdot14+\dotstwo$ The last column contains the aggregates from the 10 individual columns. The numbers 14. 15. 16 in the first vertical column are superfluous here, since no groups of a hundred contain so many primes; but they come into use on the following sheets. Finally, the 10 pages are again combined into 1, and thus cover the entire 2nd million.

But it is time to break off. \mbox{— — —} With heartfelt wishes for your well-being
\begin{center}Always yours\end{center}
\noindent Göttingen, 24 December 1849.\hfill C. F. \textsc{Gauss}.
\BookRule{.25\textwidth}
\par\endgroup

